% ============================================================
\chapter{Bayesian Hypothesis Testing}
\label{ch:testing}
% ============================================================

The frequentist hypothesis test relies on the $p$-value: the probability of observing data at least as extreme as those obtained, under the null hypothesis. But this quantity answers a question that nobody really asks: what we want to know is the probability that the hypothesis is true \emph{given the data}, not the other way round. The Bayesian test answers this directly through the \emph{Bayes factor}: the ratio of marginal likelihoods under each hypothesis. Harold Jeffreys, in his 1961 treatise, systematised this approach and proposed the scale that bears his name for interpreting the strength of evidence.

\begin{intuition}[title=\textbf{Central idea}]
Bayesian testing compares hypotheses via their posterior probabilities or
the \textbf{Bayes factor}, which measures the relative strength of
evidence for each hypothesis without depending on the stopping rule or
sample space.
\end{intuition}

% ------------------------------------------------------------
\section{Bayesian formulation of testing}
% ------------------------------------------------------------

\begin{definition}[Bayesian test]
Consider two hypotheses $H_0 : \theta \in \Theta_0$ and
$H_1 : \theta \in \Theta_1$ with $\Theta_0 \cup \Theta_1 = \Theta$
and $\Theta_0 \cap \Theta_1 = \emptyset$. Assign prior probabilities
$\PP(H_0) = \pi_0$ and $\PP(H_1) = \pi_1 = 1 - \pi_0$. The posterior
probabilities are:
\[
  \PP(H_k \mid x) = \frac{m_k(x)\,\pi_k}{m(x)},
  \quad k = 0, 1,
\]
where $m_k(x) = \int_{\Theta_k} f(x \mid \theta)\,\pi(\theta \mid H_k)\,d\theta$.
\end{definition}

\begin{definition}[Bayes factor]
The \textbf{Bayes factor} in favor of $H_0$ over $H_1$ is:
\[
  B_{01} = \frac{m_0(x)}{m_1(x)}
  = \frac{\PP(H_0 \mid x) / \PP(H_1 \mid x)}
         {\PP(H_0) / \PP(H_1)}.
\]
\end{definition}

\begin{keyformulas}[title=\textbf{Posterior odds ratio}]
\[
  \underbrace{\frac{\PP(H_0 \mid x)}{\PP(H_1 \mid x)}}_{\text{posterior odds}}
  = \underbrace{B_{01}}_{\text{Bayes factor}}
  \times
  \underbrace{\frac{\pi_0}{\pi_1}}_{\text{prior odds}}.
\]
\end{keyformulas}

% ------------------------------------------------------------
\section{Interpreting the Bayes factor}
% ------------------------------------------------------------

\begin{keyformulas}[title=\textbf{Jeffreys' scale for $B_{01}$}]
\begin{center}
\renewcommand{\arraystretch}{1.3}
\begin{tabular}{cl}
\toprule
$\log_{10} B_{01}$ & \textbf{Evidence in favor of $H_0$} \\
\midrule
$0$ to $1/2$ & Weak \\
$1/2$ to $1$ & Substantial \\
$1$ to $2$ & Strong \\
$> 2$ & Decisive \\
\bottomrule
\end{tabular}
\end{center}
\end{keyformulas}

\begin{remark}
The Bayes factor is symmetric: $B_{10} = 1/B_{01}$. If $B_{01} > 1$,
the data support $H_0$; if $B_{01} < 1$, they support $H_1$.
\end{remark}

% ------------------------------------------------------------
\section{Testing a point hypothesis}
% ------------------------------------------------------------

\begin{definition}[Point hypothesis]
We test $H_0 : \theta = \theta_0$ against
$H_1 : \theta \neq \theta_0$. We assign a point mass
$\pi_0 = \PP(\theta = \theta_0) > 0$ and a continuous density
$\pi_1(\theta)$ on $\Theta_1$.
\end{definition}

\begin{theorem}[Bayes factor for point $H_0$]
\[
  B_{01} = \frac{f(x \mid \theta_0)}
  {\int_{\Theta_1} f(x \mid \theta)\,\pi_1(\theta)\,d\theta}
  = \frac{L(\theta_0 \mid x)}{m_1(x)}.
\]
\end{theorem}

\begin{warning}[title=\textbf{Jeffreys--Lindley paradox}]
With a vague prior ($\tau_0^2 \to \infty$) under $H_1$, the Bayes
factor $B_{01} \to \infty$, always favoring $H_0$, even when the
frequentist $p$-value is very small. This paradox illustrates the
importance of the prior choice under $H_1$.
\end{warning}

\begin{example}
\textbf{Testing $\mu = 0$ in the Gaussian model.}
Let $X_1, \ldots, X_n \overset{\text{iid}}{\sim} \mathcal{N}(\mu, \sigma^2)$
with $\sigma^2$ known. Under $H_0 : \mu = 0$, under
$H_1 : \mu \sim \mathcal{N}(0, \tau^2)$.

The Bayes factor is:
\[
  B_{01} = \sqrt{\frac{\tau^2 + \sigma^2/n}{\sigma^2/n}}
  \exp\!\left(-\frac{n\bar{x}^2}{2\sigma^2}
  \cdot \frac{\tau^2}{\tau^2 + \sigma^2/n}\right).
\]
\end{example}

% ------------------------------------------------------------
\section{Bayesian decision rule}
% ------------------------------------------------------------

\begin{definition}[Optimal decision rule]
Under 0--1 loss:
\[
  L(H_k, d) = \begin{cases}
  0 & \text{if } d = k, \\
  1 & \text{if } d \neq k,
  \end{cases}
\]
the optimal decision is: accept $H_0$ if
$\PP(H_0 \mid x) > \PP(H_1 \mid x)$, i.e., $B_{01} > \pi_1/\pi_0$.

With $\pi_0 = \pi_1 = 1/2$, this reduces to accepting $H_0$ if $B_{01} > 1$.
\end{definition}

\begin{theorem}[Decision rule under asymmetric loss]
Under the loss $L(H_0, d=1) = c_0$ (type I error) and
$L(H_1, d=0) = c_1$ (type II error), the optimal rule is: accept
$H_0$ if
\[
  B_{01} > \frac{c_0}{c_1} \cdot \frac{\pi_1}{\pi_0}.
\]
\end{theorem}

% ------------------------------------------------------------
\section{Computing the Bayes factor}
% ------------------------------------------------------------

\subsection{Conjugate case}

\begin{proposition}[Bayes factor via conjugacy]
For a conjugate model, the marginal likelihood is often available
analytically. For the Bernoulli--Beta model with
$\theta \sim \text{Be}(a, b)$:
\[
  m(x) = \frac{B(a + s, b + n - s)}{B(a, b)},
\]
where $B(\cdot, \cdot)$ is the beta function.
\end{proposition}

\subsection{Savage--Dickey method}

\begin{theorem}[Savage--Dickey density ratio]
To test $H_0 : \theta = \theta_0$ against
$H_1 : \theta \sim \pi_1(\theta)$, if $\pi_1$ is the prior under $H_1$
and $\pi_1(\theta \mid x)$ the corresponding posterior:
\[
  B_{01} = \frac{\pi_1(\theta_0 \mid x)}{\pi_1(\theta_0)}.
\]
\end{theorem}

\begin{proof}
By definition:
\[
  B_{01} = \frac{f(x \mid \theta_0)}{m_1(x)}.
\]
But $\pi_1(\theta_0 \mid x) = f(x \mid \theta_0)\pi_1(\theta_0)/m_1(x)$,
so $f(x \mid \theta_0)/m_1(x) = \pi_1(\theta_0 \mid x)/\pi_1(\theta_0)$.
\end{proof}

\subsection{Numerical approximations}

\begin{definition}[Harmonic mean estimator]
The marginal likelihood can be estimated by the harmonic mean of
likelihoods evaluated on MCMC samples:
\[
  \hat{m}(x) = \left(\frac{1}{T}\sum_{t=1}^T \frac{1}{L(\theta^{(t)} \mid x)}\right)^{-1}.
\]
\end{definition}

\begin{warning}[title=\textbf{Instability of the harmonic mean}]
This estimator has potentially infinite variance and should be used
with great caution. More stable methods such as \textit{bridge sampling}
or WAIC are preferred.
\end{warning}

% ------------------------------------------------------------
\section{Bayesian information criteria}
% ------------------------------------------------------------

\begin{definition}[BIC --- Bayesian Information Criterion]
\[
  \text{BIC} = -2\log L(\hat{\theta}_{\text{MLE}} \mid x) + k\log n,
\]
where $k$ is the number of parameters. BIC approximates the log Bayes
factor: $-2\log B_{01} \approx \text{BIC}_0 - \text{BIC}_1$.
\end{definition}

\begin{definition}[WAIC --- Widely Applicable IC]
\[
  \text{WAIC} = -2\sum_{i=1}^n \log \E_{\theta \mid x}[f(x_i \mid \theta)]
  + 2\,p_{\text{WAIC}},
\]
where $p_{\text{WAIC}} = \sum_{i=1}^n \Var_{\theta \mid x}[\log f(x_i \mid \theta)]$
is the effective number of parameters.
\end{definition}

\begin{definition}[LOO-CV --- Leave-One-Out Cross-Validation]
\[
  \text{LOO} = -2\sum_{i=1}^n \log f(x_i \mid x_{-i}),
\]
where $f(x_i \mid x_{-i}) = \int f(x_i \mid \theta)\,\pi(\theta \mid x_{-i})\,d\theta$
is the leave-one-out predictive density.
\end{definition}

% ------------------------------------------------------------
\section{ROPE and practical equivalence}
% ------------------------------------------------------------

\begin{definition}[ROPE --- Region of Practical Equivalence]
The \textbf{ROPE} is an interval
$[\theta_0 - \varepsilon, \theta_0 + \varepsilon]$ defining practical
equivalence to $\theta_0$. The decision is:
\begin{itemize}
  \item Accept $H_0$ if $\text{HDI} \subset \text{ROPE}$.
  \item Reject $H_0$ if $\text{HDI} \cap \text{ROPE} = \emptyset$.
  \item Undecided otherwise.
\end{itemize}
\end{definition}

% ------------------------------------------------------------
\section{Comparison with $p$-values}
% ------------------------------------------------------------

\begin{proposition}[Bayes--$p$-value calibration]
Sellke, Bayarri, and Berger (2001) show that:
\[
  B_{01} \geq -\frac{1}{e \cdot p \cdot \log p}
  \quad \text{for } p < 1/e,
\]
where $p$ is the $p$-value. Thus, a $p$-value of $0.05$ corresponds to
$B_{01} \geq 2.5$ at most, representing weak to moderate evidence for
$H_0$.
\end{proposition}

% ------------------------------------------------------------
\section{Python implementation}
% ------------------------------------------------------------

\begin{codeblock}[title=\textbf{Bayes factor for a Gaussian test}]
\begin{minted}{python}
import numpy as np
from scipy import stats

# Data
np.random.seed(42)
n = 25
sigma = 1.0
mu_true = 0.3
data = np.random.normal(mu_true, sigma, n)
x_bar = data.mean()

# H0: mu = 0, H1: mu ~ N(0, tau^2)
tau = 1.0

# Analytic Bayes factor
var_ratio = tau**2 / (tau**2 + sigma**2/n)
log_BF01 = 0.5 * np.log(1 + n*tau**2/sigma**2) \
           - 0.5 * n * x_bar**2 / sigma**2 * var_ratio
BF01 = np.exp(log_BF01)

print(f"x_bar = {x_bar:.4f}")
print(f"BF01 = {BF01:.4f}")
print(f"log10(BF01) = {np.log10(BF01):.4f}")

if BF01 > 1:
    print("Data favor H0")
else:
    print(f"Data favor H1 (BF10 = {1/BF01:.4f})")

# Comparison with p-value
t_stat = x_bar / (sigma / np.sqrt(n))
p_value = 2 * (1 - stats.norm.cdf(abs(t_stat)))
print(f"Two-sided p-value: {p_value:.4f}")
\end{minted}
\end{codeblock}

\begin{codeblock}[title=\textbf{ROPE and HDI with PyMC}]
\begin{minted}{python}
import pymc as pm
import arviz as az
import matplotlib.pyplot as plt

# Bayesian model
with pm.Model() as model_test:
    mu = pm.Normal("mu", mu=0, sigma=10)
    y = pm.Normal("y", mu=mu, sigma=sigma,
                  observed=data)
    trace = pm.sample(3000, random_seed=42)

# HDI and ROPE
hdi = az.hdi(trace, var_names=["mu"], hdi_prob=0.95)
print("95% HDI for mu:", hdi)

# ROPE test
rope = [-0.1, 0.1]
az.plot_posterior(trace, var_names=["mu"],
                 rope=rope, ref_val=0,
                 figsize=(8, 4))
plt.savefig("rope_test.pdf")

# Proportion of posterior in ROPE
mu_samples = trace.posterior["mu"].values.flatten()
in_rope = np.mean((mu_samples >= rope[0]) &
                   (mu_samples <= rope[1]))
print(f"P(mu in ROPE | data) = {in_rope:.4f}")
\end{minted}
\end{codeblock}

\begin{codeblock}[title=\textbf{Model comparison with WAIC}]
\begin{minted}{python}
# Two competing models
with pm.Model() as model_linear:
    a = pm.Normal("a", 0, 10)
    b = pm.Normal("b", 0, 10)
    sigma_m = pm.HalfNormal("sigma", 5)
    x_pred = np.linspace(0, 10, n)
    mu_pred = a + b * x_pred
    y_obs = pm.Normal("y", mu=mu_pred, sigma=sigma_m,
                      observed=data)
    trace_lin = pm.sample(2000, random_seed=42)

with pm.Model() as model_null:
    a = pm.Normal("a", 0, 10)
    sigma_m = pm.HalfNormal("sigma", 5)
    y_obs = pm.Normal("y", mu=a, sigma=sigma_m,
                      observed=data)
    trace_null = pm.sample(2000, random_seed=42)

# WAIC comparison
comp = az.compare({"linear": trace_lin,
                    "null": trace_null})
print(comp)
\end{minted}
\end{codeblock}

% ------------------------------------------------------------
\section{Exercises}
% ------------------------------------------------------------

\begin{exercise}
Compute the Bayes factor for the test $H_0 : \theta = 1/2$ against
$H_1 : \theta \sim \text{Be}(1,1)$ in the Bernoulli model with
$n = 20$ and $s = 14$.
\end{exercise}

\begin{exercise}
Prove the Savage--Dickey density ratio. Verify numerically on the
Normal--Normal model.
\end{exercise}

\begin{exercise}
Illustrate the Jeffreys--Lindley paradox: for $n = 100$,
$\bar{x} = 0.2$, $\sigma = 1$, compute the $p$-value and the Bayes
factor $B_{01}$ for different values of $\tau \in \{1, 10, 100\}$.
Comment.
\end{exercise}

\begin{exercise}
\textbf{(Numerical)} Compare WAIC and LOO-CV for selecting between a
linear and a quadratic model on simulated data.
\end{exercise}

\begin{exercise}
Derive the BIC approximation to the Bayes factor from the Laplace
approximation of the marginal likelihood.
\end{exercise}
